\documentclass[10pt,a4paper]{amsart}
\usepackage[margin=19mm]{geometry}
\usepackage{amsmath,amssymb,mathtools}
\usepackage[scr=rsfs,cal=euler]{mathalfa}
\usepackage{xfrac}
\usepackage[unicode,hidelinks,bookmarks=false]{hyperref}
\newcommand{\ie}{i.e.\ }
\newcommand{\cf}{cf.\ }
\newcommand{\resp}{respectively\ }
\newcommand{\Subsection}{\subsection}
\providecommand{\nobreakdash}{\nobreak\hbox{-}}
\setlength{\emergencystretch}{2em}
\multlinegap0pt
\usepackage{bm}
\usepackage{bbm}
\usepackage{dsfont}
\usepackage{xspace}
\usepackage{cancel}
\usepackage{mathtools}
\usepackage{amsthm}
\makeatletter
\@ifundefined{lem}{\newtheorem{lem}{Lem}}{}
\@ifundefined{pro}{\newtheorem{pro}[lem]{Proposition}}{}
\makeatother
\DeclareMathOperator{\Hom}{Hom}
\DeclareMathOperator{\im}{Im}
\title[Some remarks on Grothendieck pairs]{Some remarks on Grothendieck pairs}
\author{Andrei Jaikin-Zapirain, Alexander Lubotzky}
\date{Source: arXiv:2401.02229v2 (CC BY 4.0)}
\begin{document}
\maketitle

\noindent\emph{Source and licence.} Some remarks on Grothendieck pairs. Version arXiv:2401.02229v2 (CC BY 4.0), distributed under the Creative Commons Attribution 4.0 International licence, as recorded in the arXiv licence metadata for this exact version and shown on the abstract page https://arxiv.org/abs/2401.02229v2. Full rights notice: licenses/arxiv-2401.02229-CC-BY-4.0.txt.\par\smallskip

\noindent\textit{Editorial scope.} Counted unit: the Pro whose source label is \texttt{gr1} in the article (source0.tex lines 303--305, source label \texttt{gr1}), delivered together with the article's own complete original proof of it (source0.tex lines 306--332, one \texttt{proof} environment of 27 lines). No mathematics is written, repaired or re-derived: statement and proof are the article's own text line for line. Required context retained uncounted: none. Every definition, display and label of those lines is retained unchanged. Omitted from this package and not redistributed: 1--302, 333--912. Nothing inside a retained excerpt is omitted or altered: every retained body is delivered exactly as the article's lines are written, so no omission receipt of any kind accompanies this package. Citations: the retained excerpts carry no citation command at all, so no bibliography entry of the article is transcribed and no reference is redistributed. Reference adaptation: none was needed and none was made; every reference inside the delivered unit resolves inside it, so no unresolved reference or citation is produced. Printed slips kept literal: no printed word, symbol, hypothesis or number of the counted statement or of its proof was altered. Layout and class: the article's own class is \texttt{amsart} with packages including amsmath, amsfonts, amssymb, amsthm, hyperref, mathtools, pdfpages, textcomp, ragged2e, blindtext, thmtools, cleveref, stmaryrd, tikz-cd; the wrapper is the lane's pinned \texttt{amsart} editorial wrapper at 10pt on A4, which restates the theorem-like environments the retained text needs and adds the article's own macro definitions that the retained text needs (\texttt{\textbackslash Hom}, \texttt{\textbackslash im}), with extra wrapper packages (bm, bbm, dsfont, xspace, cancel, mathtools). The delivered numbering is the unit's own. Assets: no retained span contains a figure environment, a graphics-inclusion command, a PDF-page inclusion, a table or a TikZ picture; no image file is copied into this package, the package declares no asset dependency and no image file is redistributed with it. Licence: version arXiv:2401.02229v2 of this article is distributed under the Creative Commons Attribution 4.0 International licence, as recorded in the arXiv licence metadata for this exact version and shown on the abstract page https://arxiv.org/abs/2401.02229v2. A full rights notice is delivered at licenses/arxiv-2401.02229-CC-BY-4.0.txt. Characters: every retained excerpt is pure ASCII, so the delivered engine, XeLaTeX with its default Latin Modern text font, reports no missing character.
   \begin{pro}\label{gr1}
Let $G\in \mathcal C$. If $H$ is commensurable with $G$, then $H\in \mathcal C$.
\end{pro}

\begin{proof}
Let  $\phi:\Lambda \hookrightarrow \Gamma$ be a Grothendieck pair. We identify the elements of $\Lambda$ with their images under the map $\phi$. Thus, we can view $\Lambda$ as a subgroup of $\Gamma$.
 We want to show that the induced map 
$$ \phi^{\#}_H:\Hom (\Gamma, H)\to \Hom (\Lambda, H), \ \tau\mapsto \tau\circ \phi,$$  is bijective. It is enough to consider two cases: $H$ is of finite index in $G$ or $G$ is of finite index in $H$.

In the first case, since $\phi_G^{\#}$ is injective,  $\phi_H^{\#}$ is injective as well. Let $\psi\in \Hom (\Lambda, H)$. Since $\phi_G^{\#}$ is surjective, there exists $\tau: \Gamma\to G$ such that $\psi=\tau\circ \phi$. Since $H$ is of finite index in $G$, $\im \tau \cap H$ is of finite index in $\im \tau$. Since $\widehat \phi$ is onto, $\im \psi$ is profinitely dense in $\im \tau$. Thus, since $\im \phi\le \im \tau \cap H$,  we obtain that
$\im \tau\le H$. Therefore, $\phi_H^{\#}$ is surjective.

Assume now that $G$ is a subgroup of $H$ of finite index. Since we have already proved the previous case, we can substitute $G$ by its core in $H$ and  assume that $G$ is also normal in $H$.

Let $\Lambda_1\le \Lambda $ ($\Gamma_1\le \Gamma$) be the intersection of all subgroups of $\Lambda$ ($\Gamma$) of index at most $|H:G|$. Then, it is clear that the restriction of $\phi$ on $\Lambda_1$, $\phi_1:\Lambda_1 \hookrightarrow \Gamma_1$ is also a Grothendieck pair. Thus $(\phi_1)_G^{\#}$ is a bijection. 

Let $\psi\in \Hom(\Lambda, H)$. Then, $\psi^{-1}(G)$ is of index at most $|H:G|$ in $\Lambda$, and so,   by our construction of $\Lambda_1$, $\Lambda_1\le \psi^{-1}(G)$. Hence $\psi(\Lambda_1)\le G$.  Consider the restriction of $\psi$ on $\Lambda_1$, $\psi_1:\Lambda_1 \to  G$.  Since $G\in \mathcal C$, there exists a unique homomorphism $\tau_1:\Gamma_1\to G$ such that $\psi_1=\tau_1\circ \phi_1$.  

If there are two homomorphisms $\tau, \tau^\prime: \Gamma\to G$ that extend $\psi$, then their restrictions on $\Gamma_1$ coincide with $\tau_1$. Since $\phi:\Lambda \hookrightarrow \Gamma$ is a Grothendieck pair, $\Gamma=\Gamma_1\Lambda$.  Thus,  $\tau=\tau^\prime$. This shows that $\phi_H$ is injective.

In order to show that $\phi_H$ is surjective, we use again that $\Gamma=\Gamma_1\Lambda$. For every  $a\in \Gamma_1$ and  every $b\in \Lambda$ we define 
$$\tau(ab)=\tau_1(a)\psi(b).$$ Since $\Gamma_1\cap \Lambda=\Lambda_1$ and $(\tau_1)_{|\Lambda_1}=\psi_1= \psi_{|\Lambda_1}$, $\tau$ is a well-defined. 
Let us show that $\tau$ is a  homomorphism.  Let  $b\in \Lambda$. Define two maps $\Gamma_1\to G$ by 
$$\alpha_1: a\mapsto \psi(b)\tau_1(a) \psi(b)^{-1}\textrm{\ and \ } \alpha_2: a\mapsto  \tau_1(bab^{-1}). $$
Their restrictions on $\Lambda_1$ coincide. Hence, since $G\in \mathcal C$, $\alpha_1=\alpha_2$. Therefore, for every  $a_1,a_2\in \Gamma_1$ and  every $b_1,b_2\in \Lambda$, we obtain that
\begin{multline*}
\tau(a_1b_1a_2b_2)=\tau(a_1b_1a_2b_1^{-1}b_1b_2)=\tau_1(a_1b_1a_2b_1^{-1})\psi(b_1b_2)=\\
\tau_1(a_1)\tau_1(b_1a_2b_1^{-1})\psi(b_1b_2)=\tau_1(a_1)\psi(b_1)\tau_1(a_2)\psi(b_1)^{-1}\psi(b_1b_2)=\\ \tau_1(a_1)\psi(b_1)\tau_1(a_2)\psi(b_2)=\tau(a_1b_1)\tau(a_2b_2).
\end{multline*}
Hence $\tau$ is a homomorphism which extends $\psi$,  and so $\phi_H$ is surjective.
\end{proof}

\end{document}
